\section{Conclusion}
\label{sec:conclusion}

Using one instrumented model and a nested protocol, we find that multimodal fusion does not improve classroom attention estimation. On a DIPSER slice, no fusion configuration exceeds face mesh with the soft target ($0.256$ against $0.124$), the gate does not differ from mean pooling, and only the input-to-head link and the soft target help; fusion also brings no gain on DAiSEE and EngageNet. On the full DIPSER release the fusion result replicates, but the per-student signal is weak and largely shared by the class, and hyper-parameter search, richer and local moments, a face crop, frozen pretrained face encoders, and temporal and class context fail to improve it. The evaluation split has a larger effect than any modelling choice: one model obtains QWK \fcLeakRandQWK{} on random windows, \fcLeakPersonQWK{} on held-out students and \fcLogoOwnQWK{} on a held-out class, with accuracy near the \fcMajority\% majority rate. Sparse probes of the raters' label raise accuracy above 70\%, but label persistence explains almost all of this gain, and self-report probes do not reproduce it. Classroom attention systems should therefore be evaluated on held-out classes, with QWK and label-only references. Future work will fine-tune face encoders and evaluate them on held-out classes.
